% !TEX TS-program = lualatex
%\documentclass[handout]{beamer}
\documentclass{beamer}
\input{../common/misomva}
\begin{document}
% Topic-specific title slide
\newcommand{\topictitle}{Introduction to Graphs in ML}
\newcommand{\topicsubtitle}{Course Overview and Motivation}
\newcommand{\misoclasstinytitleslide}{Partially based on material by:  Andreas Krause, \\[-1em] Branislav Kveton, Michael Kearns}

\MakeTopicSlide

%%%%%%%%%%%%%%%%%%%%%%%%%%%%%%%%%%%%%%%%%%%%%%%%%%%%%%%%%%%%
\begin{frame}
\frametitle{What will you learn in the Graphs in ML course?}

\textbf{Concepts\pause, tools\pause, and methods\pause} to work with graphs in ML.  \\[2em]

\pause Specific applications of graphs in ML. \\[2em]

\pause Theoretical toolbox to analyze graph-based algorithms.\\[2em]


\pause How to tackle: \emph{large graphs, online setting, graph construction} \dots  \\[2em]

\pause One example: \yellbox{Online Semi-Supervised Face Recognition} 

\end{frame}
%%%%%%%%%%%%%%%%%%%%%%%%%%%%%%%%%%%%%%%%%%%%%%%%%%%%%%%%%%%%

\begin{frame}
\frametitle{Two (main) sources of graphs in ML}

\textbf{Natural graphs as models for networks}
\begin{itemize}\addtolength{\itemsep}{.25\baselineskip}
\item<2-> {given as an input} 
\item<3-> discover interesting properties of the structure
%\item is the object of study (anomaly detection)
\end{itemize}

%\pause
\vspace{\baselineskip}
\textbf{Constructed graphs as nonparametric basis}
\begin{itemize}\addtolength{\itemsep}{0.25\baselineskip}
\item<4-> we create (learn) the similarity structure from flat data
\item<5-> it's a tool (e.g., nonparametric regularizer) to encode structural properties (e.g., independence, $\dots$)
\end{itemize}
\end{frame}


\MakeLastSlide
\end{document}
